\section{Jacobian and singularities}
\label{sec:jacobian}

\textbf{Construction from the PoE model.}
The MATLAB PoE routine keeps the intermediate transforms. These give each
joint's current unit direction $w_i$ and a point $p_i$ on its axis, both in
\texttt{base\_link}. For a revolute joint, the Jacobian column is
\begin{equation}
 {}^0J_k(:,i)=
 \begin{pmatrix}w_i\times(p_k-p_i)\\w_i\end{pmatrix},
 \qquad i=1,\ldots,6.
 \label{eq:jacobian-column}
\end{equation}
The upper rows give linear velocity and the lower rows give angular
velocity. We use $k=6$ for the endpoint. Choosing $k=5$ instead refers the
same tool motion to the wrist centre.

\textbf{Wrist Jacobian.}
Here ${}^3J_5$ denotes this full Jacobian at the wrist centre, expressed
along the DH frame-3 axes. Both velocity blocks are rotated by
$({}^0R_3)^T$; the velocities remain relative to the base.
All three wrist axes pass through $p_5$, so their linear columns vanish.
With $c_i=\cos q_i$ and $s_i=\sin q_i$, the result has the form
\begin{equation}
 {}^3J_5=\begin{pmatrix}A&0\\C&B\end{pmatrix},\qquad
 B=\begin{pmatrix}
 0&-s_4&c_4s_5\\
 0&c_4&s_4s_5\\
 1&0&c_5
 \end{pmatrix}.
 \label{eq:jacobian-wrist}
\end{equation}
The wrist subchain consists of the last three columns. Here $A$ maps the
first three joint rates to wrist-centre velocity, and $B$ maps the wrist
rates to angular velocity.

\textbf{Symbolic and numerical routes.}
The same MATLAB function accepts numerical or symbolic inputs. Numerical
inputs evaluate the Jacobian at a chosen configuration. Symbolic inputs
give an exact expression, from which the determinant is obtained.
For the full endpoint Jacobian $J={}^0J_6$, it factors as
\begin{equation}
 \det J(q)=\sin q_5\,D(q_2,q_3),\qquad D(q_2,q_3)=-\det A.
 \label{eq:jacobian-determinant}
\end{equation}
Moving the velocity reference from the wrist centre to the endpoint does
not change this determinant. The wrist factor vanishes at $q_5=0$ or
$\pi$ modulo $2\pi$. Axes 4 and 6 then lie on the same line, so one
independent rotation is lost. The arm factor depends jointly on $q_2$ and
$q_3$. Thus the wrist and arm conditions separate, while the arm condition
must be studied in the $(q_2,q_3)$ plane.

A full configuration is regular only when both factors are nonzero.
For numerical checks, the smallest singular value of the full $J$ gives
a measure of proximity to a singularity.


\textbf{How the intermediate PoE transforms give the axes.}
Write the product before joint $i$ as
$F_{i-1}=e^{\widehat\xi_1q_1}\cdots e^{\widehat\xi_{i-1}q_{i-1}}
=(R_{i-1},t_{i-1})$.
Then $w_i=R_{i-1}\omega_i$ and $p_i=R_{i-1}r_i+t_{i-1}$.
The current screw's linear component is $-w_i\times p_i$.
Evaluating its velocity field at $p_k$ adds $w_i\times p_k$, giving
the upper block of \eqref{eq:jacobian-column}.

\textbf{Why the determinant separates.}
The zero block in \eqref{eq:jacobian-wrist} gives
$\det({}^3J_5)=\det A\det B$, with $\det B=-\sin q_5$.
Expressing both velocity blocks in another rotated frame multiplies the
Jacobian by $\operatorname{diag}(R^T,R^T)$, whose determinant is one.
If $d=p_\ell-p_j$, changing the velocity reference gives
\begin{equation}
 {}^0J_\ell=\begin{pmatrix}I_3&-[d]_\times\\0&I_3\end{pmatrix}{}^0J_j.
\end{equation}
This matrix also has determinant one. Thus
$\det({}^iJ_j)=\det({}^3J_5)$ for every expression frame $i$ and velocity
reference point $p_j$, when all six columns describe the same full tool motion.
This is why ${}^3J_5$ suffices to study the full robot's singularities.
Changing the reference point preserves rank, but can change the singular
values. Numerical singular-value margins should therefore use a fixed
reference point.

\textbf{Explicit arm factor.}
Using the DH lengths and $\gamma=\pi/6$, define
\begin{equation}
 U=a_2+d_4c_3,\qquad X=Uc_2+d_3s_2,\qquad Z=-Us_2+d_3c_2.
\end{equation}
The symbolic result can then be written as
\begin{equation}
 D(q_2,q_3)=\frac{d_4}{2}
 \left[a_1Zc_3-\sqrt3\,a_1Us_3
       -\sqrt3\,X(a_2s_3-d_2c_3)\right].
 \label{eq:jacobian-arm-factor}
\end{equation}
This expression contains neither $q_1$, $q_4$ nor $q_6$.
Its zeros generally depend on the pair $(q_2,q_3)$, so the arm singularities
are studied as curves in that plane.

\textbf{From joint space to the wrist-centre workspace.}
Set $V=d_2+d_4s_3$, $c_\gamma=\cos\gamma$ and $s_\gamma=\sin\gamma$.
At $q_1=0$, the wrist centre is
\begin{equation}
 p_w=\begin{pmatrix}
 a_1+X\\c_\gamma V-s_\gamma Z\\d_1+s_\gamma V+c_\gamma Z
 \end{pmatrix},\qquad
 \rho_w=\sqrt{(a_1+X)^2+(c_\gamma V-s_\gamma Z)^2}.
\end{equation}
Changing $q_1$ rotates this point about the base $z$-axis, leaving
$(\rho_w,z_w)$ unchanged. The right panel of
Figure~\ref{fig:arm-singularity-map} maps each point of $D=0$ through
these equations. It shows the wrist centre, whose position depends on
joints 1--3. The tool endpoint also depends on the wrist orientation
through $p_E=p_w+d_6R_Ee_x$.

\textbf{Checking the cusps.}
The four marked points were checked using the arm IK polynomial
$P(t)$ in \eqref{eq:ik-quartic}. For each wrist point, form the target
invariants $h=z_w-d_1$ and $\mathcal R=x_w^2+y_w^2+h^2$.
At $t=\tan(q_3/2)$, each marked point satisfies
\begin{equation}
 P(t)=P'(t)=P''(t)=0,\qquad P'''(t)\ne0.
\end{equation}
Thus three arm IK solutions coincide at each cusp.
After normalising the polynomial coefficients, the largest residual
in the three zero conditions was $4.0\times10^{-15}$.
The cusp criterion establishes that the positioning chain is cuspidal
(see \href{https://doi.org/10.1016/j.mechmachtheory.2022.104729}
{Salunkhe et al., 2022}). This means that some arm IK solutions can be
connected without crossing an arm singularity. A motion of the full
robot must also keep $\sin q_5\ne0$ and respect the joint limits.
